% latexmk -pdf PPT12-BBDD.tex && find . -maxdepth 1 -type f ! -name "*.tex" ! -name "*.pdf" -delete

% --- CORRECCION: Forzar interpretacion raw de la entrada ---
\UseRawInputEncoding
\documentclass[aspectratio=169, 10pt, xcolor=table]{beamer}

% --- Configuracion del Tema ---
\usetheme{Madrid}
\usecolortheme{dolphin}

% --- Paquetes necesarios ---
\usepackage[utf8]{inputenc}
\usepackage[T1]{fontenc}
\usepackage[spanish,es-noquoting]{babel}
\usepackage{amsmath, amssymb}
\usepackage{mathtools}
\usepackage{booktabs}
\usepackage{graphicx}
\usepackage{listings}
\usepackage{tikz}
\usetikzlibrary{arrows.meta, positioning, shapes.geometric, babel}
\usepackage{etoolbox}
% \usepackage{upquote} % Eliminado por posibles conflictos

% --- Estilo para codigo SQL y Python ---
\definecolor{codegreen}{rgb}{0,0.6,0}
\definecolor{codegray}{rgb}{0.5,0.5,0.5}
\definecolor{codepurple}{rgb}{0.58,0,0.82}
\definecolor{backcolour}{rgb}{0.95,0.95,0.92}
\lstdefinestyle{mysql}{
    language=SQL,
    backgroundcolor=\color{backcolour},
    commentstyle=\color{codegreen},
    keywordstyle=\color{blue},
    numberstyle=\tiny\color{codegray},
    stringstyle=\color{codepurple},
    basicstyle=\ttfamily\scriptsize,
    breaklines=true,
    frame=single,
    showstringspaces=false,
    literate={á}{{\'a}}1 {é}{{\'e}}1 {í}{{\'i}}1 {ó}{{\'o}}1 {ú}{{\'u}}1 {ñ}{{\~n}}1 {ü}{{\"u}}1
}
\lstdefinestyle{python}{
    language=Python,
    backgroundcolor=\color{backcolour},
    commentstyle=\color{codegreen},
    keywordstyle=\color{blue},
    numberstyle=\tiny\color{codegray},
    stringstyle=\color{codepurple},
    basicstyle=\ttfamily\scriptsize,
    breaklines=true,
    frame=single,
    showstringspaces=false,
    literate={á}{{\'a}}1 {é}{{\'e}}1 {í}{{\'i}}1 {ó}{{\'o}}1 {ú}{{\'u}}1 {ñ}{{\~n}}1 {ü}{{\"u}}1
}
\lstset{style=python} % Por defecto Python

% --- Pie de pagina personalizado ---
\makeatletter
\setbeamertemplate{footline}{
  \leavevmode%
  \hbox{%
  \begin{beamercolorbox}[wd=.333333\paperwidth,ht=2.25ex,dp=1ex,center]{author in head/foot}%
    \usebeamerfont{author in head/foot}\insertshortauthor
  \end{beamercolorbox}%
  \begin{beamercolorbox}[wd=.333333\paperwidth,ht=2.25ex,dp=1ex,center]{title in head/foot}%
    \usebeamerfont{title in head/foot}Sesión 12
  \end{beamercolorbox}%
  \begin{beamercolorbox}[wd=.333333\paperwidth,ht=2.25ex,dp=1ex,right]{date in head/foot}%
    \usebeamerfont{date in head/foot}BBDD \hfill \insertframenumber/\inserttotalframenumber \hspace*{0.3cm}
  \end{beamercolorbox}}%
  \vskip0pt%
}
\makeatother

% --- Datos de la portada ---
\title[SESIÓN 12 BBDD]{\textbf{Sesión 12}}
\subtitle{Bases de Datos Vectoriales y RAG}
\author[Prof. C. Sánchez]{Carlos César Sánchez Coronel}
\institute{\textbf{BASES DE DATOS}}
\date{2026}

\setbeamertemplate{navigation symbols}{}

\AtBeginSection[]{
  \begin{frame}
    \frametitle{Contenido de la sección}
    \tableofcontents[currentsection]
  \end{frame}
}

\begin{document}

% Desactivar caracteres activos de babel (<, >) para que no interfieran con listings
\shorthandoff{<>}

\begin{frame}
    \titlepage
\end{frame}

\begin{frame}{Contenido general}
    \tableofcontents
\end{frame}

% ============================================================================
% SECCIÓN 1: INTRODUCCIÓN A LAS BASES DE DATOS VECTORIALES
% ============================================================================
\section{Introducción a las Bases de Datos Vectoriales}

\begin{frame}{¿Qué es un vector en el contexto de IA?}
    \begin{itemize}
        \item Un \textbf{embedding} es una representación densa de baja dimensionalidad de un objeto (texto, imagen, audio, etc.) en un espacio continuo.
        \item Objetos semánticamente similares tienen vectores cercanos según alguna métrica.
        \item Generados por modelos como Word2Vec, BERT, Sentence Transformers, CLIP.
    \end{itemize}
\end{frame}

\begin{frame}{Limitaciones de las bases de datos tradicionales}
    \begin{itemize}
        \item No están optimizadas para búsquedas por similitud semántica.
        \item Búsquedas exactas o por palabras clave (LIKE, full-text) no capturan significado.
        \item No manejan eficientemente vectores de alta dimensionalidad.
        \item Escaneo completo (O(n)) inviable para millones de vectores.
    \end{itemize}
\end{frame}

\begin{frame}{¿Qué hace una base de datos vectorial?}
    \begin{itemize}
        \item Almacena vectores junto con metadatos.
        \item Indexa los vectores para búsquedas rápidas de vecinos cercanos (ANN).
        \item Ofrece APIs para calcular distancias y recuperar elementos similares.
    \end{itemize}
\end{frame}

% ============================================================================
% SECCIÓN 2: EMBEDDINGS
% ============================================================================
\section{Embeddings}

\begin{frame}{Definición y generación}
    \begin{itemize}
        \item Un \textbf{embedding} es una lista ordenada de números que representa un objeto.
        \item Modelos populares:
        \begin{itemize}
            \item Texto: Sentence Transformers (all-MiniLM-L6-v2), OpenAI embeddings.
            \item Imágenes: CLIP, ResNet.
            \item Multimodal: CLIP, ImageBind.
        \end{itemize}
    \end{itemize}
\end{frame}

\begin{frame}[fragile]{Ejemplo con Sentence Transformers}
    \begin{lstlisting}[style=python]
from sentence_transformers import SentenceTransformer

model = SentenceTransformer('all-MiniLM-L6-v2')
oraciones = ["El gato duerme en el sofá", "Un perro juega en el parque"]
embeddings = model.encode(oraciones)
print(embeddings.shape)  # (2, 384)
print(embeddings[0][:5]) # primeros 5 valores del primer vector
    \end{lstlisting}
\end{frame}

% ============================================================================
% SECCIÓN 3: BÚSQUEDA DE SIMILITUD
% ============================================================================
\section{Búsqueda de Similitud: k-NN y ANN}

\begin{frame}{k-NN exacto vs ANN}
    \begin{description}
        \item[k-NN exacto] Calcula distancia con todos los vectores, O(n) por consulta. Inviable para millones.
        \item[ANN (Approximate Nearest Neighbors)] Algoritmos que sacrifican precisión por velocidad (complejidad sublineal).
    \end{description}
\end{frame}

\begin{frame}{Algoritmos ANN}
    \begin{description}
        \item[HNSW] Grafos multinivel, alta velocidad y precisión, alto consumo de memoria.
        \item[IVF] Divide espacio en clusters (k-means), explora clusters cercanos.
        \item[PQ (Product Quantization)] Comprime vectores, reduce memoria, usado con IVF.
    \end{description}
\end{frame}

\begin{frame}{Comparativa de algoritmos}
    \begin{table}
        \centering
        \begin{tabular}{l|c|c}
            \toprule
            Algoritmo & Velocidad & Precisión \\
            \midrule
            HNSW & Alta & Muy alta \\
            IVF & Alta & Alta \\
            PQ & Muy alta & Media \\
            \bottomrule
        \end{tabular}
    \end{table}
\end{frame}

% ============================================================================
% SECCIÓN 4: MÉTRICAS DE DISTANCIA
% ============================================================================
\section{Métricas de Distancia}

\begin{frame}{Métricas comunes}
    \begin{description}
        \item[Similitud de coseno] Mide ángulo entre vectores; ignora magnitudes. $\frac{u \cdot v}{\|u\|\|v\|}$
        \item[Distancia Euclidiana (L2)] Distancia geométrica: $\sqrt{\sum (u_i - v_i)^2}$
        \item[Producto punto] Si vectores normalizados, equivale al coseno.
    \end{description}
    \begin{alertblock}{Importante}
        Usar la métrica con la que fue entrenado el modelo (ej. coseno para Sentence Transformers).
    \end{alertblock}
\end{frame}

% ============================================================================
% SECCIÓN 5: TECNOLOGÍAS DE BASES DE DATOS VECTORIALES
% ============================================================================
\section{Tecnologías de Bases de Datos Vectoriales}

\begin{frame}{Bases de datos nativas}
    \begin{description}
        \item[Pinecone] Servicio cloud gestionado, índice HNSW.
        \item[Milvus] Open source, múltiples índices, despliegue on-premise/cloud.
        \item[Weaviate] Open source, módulos de ML, búsqueda híbrida.
        \item[Qdrant] Alto rendimiento en Rust, open source y cloud.
        \item[Chroma] Ligera, para prototipos y proyectos pequeños.
    \end{description}
\end{frame}

\begin{frame}{Extensiones sobre bases de datos existentes}
    \begin{description}
        \item[pgvector] Extensión PostgreSQL, índices IVFFlat y HNSW.
        \item[Elasticsearch] Tipo dense\_vector, búsqueda híbrida.
        \item[Redis Stack] Módulo RedisVL para búsqueda vectorial.
        \item[MongoDB Atlas Search] Soporte vectorial (beta).
    \end{description}
\end{frame}

\begin{frame}{Tabla comparativa}
    \begin{table}
        \centering
        \begin{tabular}{l|c|c}
            \toprule
            Tecnología & Tipo & Índices \\
            \midrule
            Pinecone & Nativa (cloud) & HNSW \\
            Milvus & Nativa (open source) & HNSW, IVF, etc. \\
            Weaviate & Nativa (open source) & HNSW \\
            pgvector & Extensión PostgreSQL & IVFFlat, HNSW \\
            Elasticsearch & Motor de búsqueda & HNSW (plugin) \\
            RedisVL & Módulo Redis & HNSW \\
            \bottomrule
        \end{tabular}
    \end{table}
\end{frame}

\begin{frame}[fragile]{Ejemplo con pgvector}
    \begin{lstlisting}[style=mysql]
CREATE EXTENSION vector;

CREATE TABLE productos (
    id SERIAL PRIMARY KEY,
    nombre TEXT,
    descripcion TEXT,
    embedding vector(384)
);

-- Indice HNSW (requiere pgvector >= 0.5.0)
CREATE INDEX ON productos USING hnsw (embedding vector_cosine_ops);

-- Búsqueda: productos similares a una descripción
SELECT nombre, embedding <=> '[0.1,0.2,...]' AS distancia
FROM productos
ORDER BY embedding <=> '[0.1,0.2,...]'
LIMIT 10;
    \end{lstlisting}
\end{frame}

% ============================================================================
% SECCIÓN 6: RETRIEVAL AUGMENTED GENERATION (RAG)
% ============================================================================
\section{Retrieval Augmented Generation (RAG)}

\begin{frame}{Problema que resuelve RAG}
    \begin{itemize}
        \item Los LLMs tienen conocimiento limitado a su fecha de entrenamiento.
        \item Pueden alucinar (inventar) cuando no saben algo.
        \item RAG proporciona contexto relevante extraído de una base de conocimiento, reduciendo alucinaciones.
    \end{itemize}
\end{frame}

\begin{frame}{Flujo de trabajo típico}
    \begin{enumerate}
        \item \textbf{Indexación}: dividir documentos, generar embeddings, almacenar en BD vectorial.
        \item \textbf{Consulta}: usuario hace pregunta.
        \item \textbf{Embedding de consulta} y búsqueda de fragmentos similares.
        \item \textbf{Generación de prompt} con fragmentos recuperados.
        \item \textbf{LLM} genera respuesta basada en el contexto.
    \end{enumerate}
\end{frame}

\begin{frame}{Orquestadores para RAG}
    \begin{description}
        \item[LangChain] Framework para cadenas con LLMs, integración con múltiples fuentes y bases vectoriales.
        \item[LlamaIndex] Especializado en indexación y recuperación de datos.
        \item[Haystack] Framework modular para búsqueda y RAG.
    \end{description}
\end{frame}

\begin{frame}[fragile]{Ejemplo con LangChain y Chroma}
    \begin{lstlisting}[style=python]
from langchain.document_loaders import PyPDFLoader
from langchain.text_splitter import RecursiveCharacterTextSplitter
from langchain.embeddings import HuggingFaceEmbeddings
from langchain.vectorstores import Chroma
from langchain.llms import OpenAI
from langchain.chains import RetrievalQA

# Cargar y dividir PDF
loader = PyPDFLoader("documento.pdf")
documentos = loader.load()
splitter = RecursiveCharacterTextSplitter(chunk_size=500, chunk_overlap=50)
docs = splitter.split_documents(documentos)

# Embeddings y vectorstore
embeddings = HuggingFaceEmbeddings(model_name="all-MiniLM-L6-v2")
vectorstore = Chroma.from_documents(docs, embeddings)

# Cadena QA
llm = OpenAI(temperature=0)
qa = RetrievalQA.from_chain_type(llm=llm, retriever=vectorstore.as_retriever())
respuesta = qa.run("¿Cuál es el tema principal del documento?")
print(respuesta)
    \end{lstlisting}
\end{frame}

% ============================================================================
% SECCIÓN 7: INTEGRACIÓN CON SQL Y NOSQL
% ============================================================================
\section{Integración con SQL y NoSQL}

\begin{frame}{Casos de integración}
    \begin{itemize}
        \item Metadatos en PostgreSQL, embeddings en pgvector (búsqueda híbrida).
        \item Documentos en MongoDB, búsqueda semántica con base vectorial externa.
        \item Data lake con archivos Parquet + catálogo de embeddings.
    \end{itemize}
\end{frame}

\begin{frame}{Estrategias de consistencia}
    \begin{itemize}
        \item \textbf{Sincronización síncrona}: escribir en ambas bases dentro de transacción.
        \item \textbf{CDC (Change Data Capture)}: capturar cambios y propagar vía Kafka/Debezium.
        \item \textbf{Batch periódico}: reindexar cada cierto tiempo.
    \end{itemize}
\end{frame}

\begin{frame}[fragile]{Búsqueda híbrida con pgvector}
    \begin{lstlisting}[style=mysql]
SELECT nombre, precio, embedding <=> query_embedding AS distancia
FROM productos
WHERE categoria = 'electronica' AND precio < 500
ORDER BY embedding <=> query_embedding
LIMIT 10;
    \end{lstlisting}
\end{frame}

% ============================================================================
% SECCIÓN 8: DESPLIEGUE ON-PREMISE Y CLOUD
% ============================================================================
\section{Despliegue On-Premise y Cloud}

\begin{frame}{On-Premise vs Cloud}
    \begin{description}
        \item[On-Premise] Control total, seguridad, sin dependencia de proveedor. Mantenimiento y escalado manual.
        \item[Cloud] Escalabilidad elástica, menor overhead operativo, pago por uso. Dependencia del proveedor, costos variables.
    \end{description}
\end{frame}

\begin{frame}{Servicios cloud gestionados}
    \begin{itemize}
        \item \textbf{AWS}: OpenSearch Serverless (vectorial), Aurora PostgreSQL con pgvector, MemoryDB (Redis).
        \item \textbf{GCP}: Vertex AI Vector Search, AlloyDB con pgvector.
        \item \textbf{Azure}: Cognitive Search con vectorial, Cosmos DB (preview).
    \end{itemize}
\end{frame}

\begin{frame}{Dimensionamiento y costos}
    \begin{itemize}
        \item Factores: número de vectores, dimensionalidad, tipo de índice, QPS.
        \item Ejemplo Pinecone: $0.50/hora por cada millón de vectores (aprox). Para 10M vectores, costo mensual ≈ $3,600.
        \item pgvector en RDS puede ser más económico si ya se tiene instancia.
    \end{itemize}
\end{frame}

% ============================================================================
% SECCIÓN 9: CASOS DE USO REALES
% ============================================================================
\section{Casos de Uso Reales}

\begin{frame}{Búsqueda semántica en e-commerce}
    Usuarios buscan con lenguaje natural (ej. ``vestido rojo elegante''). Embeddings de descripciones vs. embedding de consulta.
\end{frame}

\begin{frame}{Sistemas de preguntas y respuestas (FAQ)}
    Documentación interna indexada en base vectorial; chatbot RAG responde preguntas de empleados.
\end{frame}

\begin{frame}{Recomendación de contenidos}
    Embeddings de películas basados en trama/género; recomendar las más similares a las vistas por el usuario.
\end{frame}

\begin{frame}{Detección de plagio / documentos duplicados}
    Comparar embeddings de documentos para encontrar similitudes.
\end{frame}

\begin{frame}{Búsqueda multimodal}
    Usar CLIP para buscar imágenes con texto o viceversa.
\end{frame}

% ============================================================================
% SECCIÓN 10: EJERCICIOS RESUELTOS
% ============================================================================
\section{Ejercicios Resueltos}

\begin{frame}[fragile]{Ejercicio 1: Embeddings y pgvector}
    \textbf{Enunciado:} Generar embeddings de oraciones y buscar similitud con pgvector.
    \begin{lstlisting}[style=python]
import psycopg2
from sentence_transformers import SentenceTransformer

conn = psycopg2.connect(dbname="test", user="postgres")
cur = conn.cursor()
cur.execute("CREATE EXTENSION IF NOT EXISTS vector;")
cur.execute("DROP TABLE IF EXISTS oraciones;")
cur.execute("CREATE TABLE oraciones (id SERIAL PRIMARY KEY, texto TEXT, embedding vector(384));")

model = SentenceTransformer('all-MiniLM-L6-v2')
oraciones = ["El gato duerme en el sofá", "Los perros juegan en el parque", "Me gusta la comida italiana", "El cielo está nublado hoy"]
embeddings = model.encode(oraciones).tolist()

for texto, emb in zip(oraciones, embeddings):
    cur.execute("INSERT INTO oraciones (texto, embedding) VALUES (%s, %s);", (texto, emb))
conn.commit()

consulta = "animales domésticos"
emb_consulta = model.encode([consulta]).tolist()[0]
cur.execute("""SELECT texto, embedding <=> %s::vector AS distancia FROM oraciones ORDER BY embedding <=> %s::vector LIMIT 2;""", (emb_consulta, emb_consulta))
for texto, dist in cur.fetchall():
    print(f"{texto} (distancia: {dist:.4f})")
cur.close(); conn.close()
    \end{lstlisting}
\end{frame}

\begin{frame}[fragile]{Ejercicio 2: Pipeline RAG con LangChain y Chroma}
    \textbf{Enunciado:} Cargar PDF, dividir, crear embeddings, construir QA con Chroma.
    \begin{lstlisting}[style=python]
from langchain.document_loaders import PyPDFLoader
from langchain.text_splitter import RecursiveCharacterTextSplitter
from langchain.embeddings import HuggingFaceEmbeddings
from langchain.vectorstores import Chroma
from langchain.llms import OpenAI
from langchain.chains import RetrievalQA

loader = PyPDFLoader("documento.pdf")
documentos = loader.load()
splitter = RecursiveCharacterTextSplitter(chunk_size=500, chunk_overlap=50)
docs = splitter.split_documents(documentos)

embeddings = HuggingFaceEmbeddings(model_name="all-MiniLM-L6-v2")
vectorstore = Chroma.from_documents(docs, embeddings)

llm = OpenAI(temperature=0)
qa = RetrievalQA.from_chain_type(llm=llm, retriever=vectorstore.as_retriever())
respuesta = qa.run("¿Cual es el tema principal del documento?")
print(respuesta)
    \end{lstlisting}
\end{frame}

\begin{frame}[fragile]{Ejercicio 3: Búsqueda híbrida con pgvector}
    \textbf{Enunciado:} Filtrar por precio y categoría, ordenar por similitud.
    \begin{lstlisting}[style=python]
import psycopg2
from sentence_transformers import SentenceTransformer

conn = psycopg2.connect(dbname="test", user="postgres")
cur = conn.cursor()

model = SentenceTransformer('all-MiniLM-L6-v2')
consulta = "teléfono inteligente con buena cámara"
emb_consulta = model.encode([consulta]).tolist()[0]

cur.execute("""
    SELECT nombre, precio, embedding <=> %s::vector AS dist
    FROM productos
    WHERE precio > 50 AND categoria = 'electronica'
    ORDER BY embedding <=> %s::vector
    LIMIT 5;
""", (emb_consulta, emb_consulta))

for row in cur.fetchall():
    print(row)
cur.close(); conn.close()
    \end{lstlisting}
\end{frame}

\begin{frame}[fragile]{Ejercicio 4: Comparación de índices en Milvus}
    \textbf{Enunciado:} Crear colección, insertar vectores, comparar HNSW vs IVFFlat.
    \begin{lstlisting}[style=python]
from pymilvus import connections, Collection, FieldSchema, CollectionSchema, DataType, utility
import numpy as np
import time

connections.connect("default", host="localhost", port="19530")

fields = [
    FieldSchema(name="id", dtype=DataType.INT64, is_primary=True),
    FieldSchema(name="embedding", dtype=DataType.FLOAT_VECTOR, dim=128)
]
schema = CollectionSchema(fields)
collection = Collection("test_vector", schema)

data = [
    [i for i in range(10000)],
    [np.random.rand(128).tolist() for _ in range(10000)]
]
collection.insert(data)
collection.flush()

# HNSW
hnsw_params = {"index_type": "HNSW", "metric_type": "COSINE", "params": {"M": 16, "efConstruction": 200}}
collection.create_index("embedding", hnsw_params)
collection.load()
start = time.time()
collection.search([np.random.rand(128).tolist()], "embedding", param={"metric_type": "COSINE", "params": {"ef": 64}}, limit=10)
print("HNSW tiempo:", time.time() - start)

# IVF (drop index first)
collection.drop_index()
ivf_params = {"index_type": "IVF_FLAT", "metric_type": "COSINE", "params": {"nlist": 128}}
collection.create_index("embedding", ivf_params)
collection.load()
start = time.time()
collection.search([np.random.rand(128).tolist()], "embedding", param={"metric_type": "COSINE", "params": {"nprobe": 10}}, limit=10)
print("IVF tiempo:", time.time() - start)
    \end{lstlisting}
\end{frame}

% ============================================================================
% SECCIÓN 11: GLOSARIO
% ============================================================================
\section{Glosario}

\begin{frame}{Términos Clave}
    \begin{description}
        \item[Embedding] Representación vectorial de un objeto.
        \item[Vector] Lista ordenada de números.
        \item[Dimensionalidad] Número de componentes del vector.
        \item[Similitud de coseno] Medida del ángulo entre vectores.
        \item[Distancia Euclidiana] Distancia geométrica.
        \item[ANN] Approximate Nearest Neighbors.
        \item[HNSW] Hierarchical Navigable Small World.
        \item[IVF] Inverted File Index.
        \item[PQ] Product Quantization.
        \item[RAG] Retrieval Augmented Generation.
        \item[LangChain] Framework para aplicaciones con LLMs.
        \item[pgvector] Extensión vectorial para PostgreSQL.
        \item[Pinecone] Base de datos vectorial gestionada.
        \item[Milvus] Plataforma de búsqueda vectorial open source.
    \end{description}
\end{frame}

% --- Diapositiva final ---
\begin{frame}
    \centering
    \Huge \textbf{¿Preguntas?}
    
    \vspace{1cm}
    \large ¡Gracias por su atención!
\end{frame}

\end{document}